\def\EMtitlesize{26}\def\EMtitleskip{31}
\begin{frame}{بخش ۳}
\begin{tikzpicture}[remember picture,overlay]
\node[fill=SLbar,draw=SLbarEdge,line width=0.8bp,rounded corners=4bp,minimum width=520bp,minimum height=170bp] (b) at ([yshift=-20bp]current page.center) {};
\node[text=SLthm,font=\bfseries\fontsize{30}{36}\selectfont] at ([yshift=36bp]b.center) {\EMdir{بخش ۳}};
\node[text=SLdef,font=\fontsize{24}{30}\selectfont] at ([yshift=-30bp]b.center) {\EMdir{آنالیز مختلط}};
\end{tikzpicture}
\end{frame}

\section{اعداد مختلط}
\def\EMfootchapter{فصل ۸ \ \lr{\textbullet}\ اعداد مختلط}
\def\EMtitlesize{26}\def\EMtitleskip{31}
\begin{frame}{فصل ۸: اعداد مختلط}
{\fontsize{20}{26}\selectfont\bfseries\color{SLdef}\EMalign مطالب این فصل\par}
\vspace{10bp}
\EMitem{\lr{\textbullet}}{\dimexpr\linewidth-26bp\relax}{مجموعۀ اعداد مختلط}
\EMitem{\lr{\textbullet}}{\dimexpr\linewidth-26bp\relax}{چهار عمل اصلی}
\EMitem{\lr{\textbullet}}{\dimexpr\linewidth-26bp\relax}{خواص جمع و ضرب}
\EMitem{\lr{\textbullet}}{\dimexpr\linewidth-26bp\relax}{مزدوج مختلط}
\EMitem{\lr{\textbullet}}{\dimexpr\linewidth-26bp\relax}{نامساوی مثلثی}
\EMitem{\lr{\textbullet}}{\dimexpr\linewidth-26bp\relax}{فرم قطبی اعداد مختلط}
\EMitem{\lr{\textbullet}}{\dimexpr\linewidth-26bp\relax}{ضرب دو عدد مختلط در فرم قطبی}
\EMitem{\lr{\textbullet}}{\dimexpr\linewidth-26bp\relax}{ریشۀ \EMim{m}-ام یک عدد مختلط}
\end{frame}

\def\EMtitlesize{26}\def\EMtitleskip{31}
\begin{frame}{مجموعۀ اعداد مختلط}
\begin{minipage}[t]{500bp}\vspace{0pt}
\begin{EMbox}{def}{1}{}
\begin{uncoverenv}<2->
\EMp{مجموعۀ اعداد مختلط تعمیم اعداد واقع بر محور اعداد حقیقی به صفحه‌ای از اعداد دارای دو مولفۀ حقیقی و موهومی می‌باشد.}
\vspace{3bp}
\EMeq{\mathbb{R}\to\mathbb{C}}
\end{uncoverenv}
\end{EMbox}
\begin{uncoverenv}<3->
\EMeq{\mathbb{C} =\{(x,y)\mid x,y\in\mathbb{R}\},\qquad z =(x,y),\qquad x \\[4pt] =\operatorname{Re}\{z\},\qquad y =\operatorname{Im}\{z\}}
\end{uncoverenv}
\EMgap
\end{minipage}
\EMside{190.0bp}{\uncover<4->{\EMfigure{1.500}{m08-complex-plane}{نمایش عدد مختلط \EMim{z=(x,y)} در صفحۀ مختلط؛ اندازۀ \EMim{|z|} و آرگومان \EMim{\theta}}}}
\end{frame}

\def\EMtitlesize{26}\def\EMtitleskip{31}
\begin{frame}{مجموعۀ اعداد مختلط (ادامه)}
\begin{uncoverenv}<2->
\EMp{اندازه یا نرم (\lr{norm}) عدد مختلط:}
\vspace{3bp}
\EMeq{|z| =\sqrt{x^2+y^2}}
\end{uncoverenv}
\EMgap
\begin{uncoverenv}<3->
\EMp{فاز یا آرگومان:}
\vspace{3bp}
\EMeq{\theta =\measuredangle z =\tan^{-1}\!\left(\frac{y}{x}\right),\qquad z =|z|\measuredangle\theta}
\end{uncoverenv}
\EMgap
\begin{uncoverenv}<4->
\EMeq{x =|z|\cos\theta,\qquad y =|z|\sin\theta}
\end{uncoverenv}
\EMgap
\begin{uncoverenv}<5->
\EMeq{z =x\times(1,0)+y\times(0,1) =x+iy}
\end{uncoverenv}
\EMgap
\begin{EMbox}{imp}{5}{}
\begin{uncoverenv}<6->
\EMeq{i^2 =-1}
\end{uncoverenv}
\end{EMbox}
\end{frame}

\def\EMtitlesize{26}\def\EMtitleskip{31}
\begin{frame}{چهار عمل اصلی}
\begin{uncoverenv}<2->
\EMp{فرض کنید \EMim{z_1=(a_1,b_1)=a_1+ib_1} و \EMim{z_2=(a_2,b_2)=a_2+ib_2}.}
\end{uncoverenv}
\EMgap
\begin{uncoverenv}<3->
\EMp{\textbf{(\lr{I}) جمع:}}
\vspace{3bp}
\EMeq{z_1+z_2 =(a_1+a_2,\,b_1+b_2) =a_1+a_2+i(b_1+b_2)}
\end{uncoverenv}
\EMgap
\begin{uncoverenv}<4->
\EMp{\textbf{(\lr{II}) تفریق:}}
\vspace{3bp}
\EMeq{z_1-z_2 =(a_1-a_2,\,b_1-b_2) =a_1-a_2+i(b_1-b_2)}
\end{uncoverenv}
\EMgap
\begin{uncoverenv}<5->
\EMp{\textbf{(\lr{III}) ضرب:}}
\vspace{3bp}
\EMeq{\begin{aligned}
z_1\times z_2&=(a_1+ib_1)\times(a_2+ib_2)\\
&=a_1\times a_2+a_1\times ib_2+ib_1\times a_2+ib_1\times ib_2\\
&=(a_1a_2-b_1b_2)+i(a_1b_2+b_1a_2)
\end{aligned}}
\end{uncoverenv}
\end{frame}

\def\EMtitlesize{26}\def\EMtitleskip{31}
\begin{frame}{چهار عمل اصلی (ادامه)}
\begin{uncoverenv}<2->
\EMp{ضرب در یک عدد حقیقی:}
\vspace{3bp}
\EMeq{\forall\alpha\in\mathbb{R}:\quad \alpha z_1 =\alpha a_1+i\alpha b_1}
\end{uncoverenv}
\EMgap
\begin{uncoverenv}<3->
\EMp{\textbf{(\lr{IV}) تقسیم:}}
\vspace{3bp}
\EMeq{\begin{aligned}
z_1\div z_2&=(a_1+ib_1)\div(a_2+ib_2)=\frac{a_1+ib_1}{a_2+ib_2}=\frac{(a_1+ib_1)(a_2-ib_2)}{(a_2+ib_2)(a_2-ib_2)}\\
&=\frac{(a_1a_2+b_1b_2)+i(b_1a_2-a_1b_2)}{a_2^2+b_2^2}
=\frac{a_1a_2+b_1b_2}{a_2^2+b_2^2}+i\,\frac{b_1a_2-a_1b_2}{a_2^2+b_2^2}
\end{aligned}}
\end{uncoverenv}
\end{frame}

\def\EMtitlesize{26}\def\EMtitleskip{31}
\begin{frame}{خواص جمع و ضرب}
\begin{EMbox}{thm}{1}{}
\begin{uncoverenv}<2->
\EMp{مجموعۀ اعداد مختلط همراه با دو عمل جمع و ضرب تشکیل یک میدان غیرترتیب‌پذیر می‌دهد.}
\end{uncoverenv}
\end{EMbox}
\end{frame}

\def\EMtitlesize{26}\def\EMtitleskip{31}
\begin{frame}{خواص جمع و ضرب (ادامه)}
\begin{uncoverenv}<2->
\EMp{برای \EMim{z_1,z_2,z_3\in\mathbb{C}}:}
\end{uncoverenv}
\EMgap
\begin{uncoverenv}<3->
\EMeq{z_1+z_2\in\mathbb{C},\qquad z_1\times z_2\in\mathbb{C}}
\end{uncoverenv}
\EMgap
\begin{uncoverenv}<4->
\EMeq{z_1+z_2 =z_2+z_1,\qquad z_1\times z_2 =z_2\times z_1}
\end{uncoverenv}
\EMgap
\begin{uncoverenv}<5->
\EMeq{(z_1+z_2)+z_3 =z_1+(z_2+z_3),\qquad (z_1z_2)z_3 =z_1(z_2z_3)}
\end{uncoverenv}
\EMgap
\begin{uncoverenv}<6->
\EMeq{z_1(z_2+z_3) =z_1z_2+z_1z_3}
\end{uncoverenv}
\EMgap
\begin{uncoverenv}<7->
\EMeq{z_1+(0,0) =z_1+0 =z_1,\qquad z_1+(-z_1) =(0,0) =0}
\end{uncoverenv}
\EMgap
\begin{uncoverenv}<8->
\EMeq{-z_1 =-(a_1,b_1) =(-a_1,-b_1) =-a_1-ib_1}
\end{uncoverenv}
\EMgap
\begin{uncoverenv}<9->
\EMeq{z_1\times(1,0) =z_1\times1 =z_1,\qquad z_1\times z_1^{-1} =(1,0) =1}
\end{uncoverenv}
\EMgap
\begin{uncoverenv}<10->
\EMeq{z_1^{-1} =\left(\frac{a_1}{a_1^2+b_1^2},\frac{-b_1}{a_1^2+b_1^2}\right) =\frac{a_1}{a_1^2+b_1^2}-i\,\frac{b_1}{a_1^2+b_1^2}}
\end{uncoverenv}
\end{frame}

\def\EMtitlesize{26}\def\EMtitleskip{31}
\begin{frame}{خواص جمع و ضرب (ادامه)}
\begin{EMbox}{rem}{1}{}
\begin{uncoverenv}<2->
\EMp{در مجموعۀ اعداد مختلط نمی‌توان رابطۀ ترتیب \EMim{z_1<z_2} تعریف کرد.}
\end{uncoverenv}
\end{EMbox}
\end{frame}

\def\EMtitlesize{26}\def\EMtitleskip{31}
\begin{frame}{مزدوج مختلط}
\begin{EMbox}{def}{1}{}
\begin{uncoverenv}<2->
\EMeq{z =x+iy\quad \Longrightarrow\quad \bar z\triangleq x-iy}
\end{uncoverenv}
\end{EMbox}
\end{frame}

\def\EMtitlesize{26}\def\EMtitleskip{31}
\begin{frame}{مزدوج مختلط (ادامه)}
\begin{minipage}[t]{500bp}\vspace{0pt}
\begin{EMbox}{ex}{1}{}
\EMp{مزدوج عدد مختلط \EMim{z=5+2i} را در صفحۀ مختلط رسم کنید.}
\end{EMbox}
\end{minipage}
\EMside{53.4bp}{\uncover<2->{\EMfigure{1.500}{m08-conjugate}{عدد مختلط \EMim{z=5+2i} و مزدوج آن \EMim{\bar z=5-2i}}}}
\end{frame}

\def\EMtitlesize{26}\def\EMtitleskip{31}
\begin{frame}{مزدوج مختلط (ادامه)}
\begin{EMbox}{exe}{1}{}
\EMp{عبارت‌های زیر را ثابت کنید.}
\vspace{3bp}
\EMeq{\operatorname{Re}\{z\} =\frac12\,(z+\bar z),\qquad \operatorname{Im}\{z\} =\frac{1}{2i}\,(z-\bar z),\qquad |z|^2 =z\cdot\bar z}
\EMgap
\begin{uncoverenv}<2->
\EMeq{\overline{z_1+z_2} =\bar z_1+\bar z_2,\qquad \overline{z_1\cdot z_2} =\bar z_1\cdot\bar z_2,\qquad \overline{\left(\frac{z_1}{z_2}\right)} =\frac{\bar z_1}{\bar z_2}}
\end{uncoverenv}
\EMgap
\begin{uncoverenv}<3->
\EMeq{|z| =|\bar z|,\qquad |\operatorname{Re}\{z\}|\le|z|,\qquad |\operatorname{Im}\{z\}|\le|z|}
\end{uncoverenv}
\end{EMbox}
\end{frame}

\def\EMtitlesize{26}\def\EMtitleskip{31}
\begin{frame}{نامساوی مثلثی}
\begin{minipage}[t]{500bp}\vspace{0pt}
\EMgap
\begin{EMbox}{exe}{1}{}
\EMitem{\lr{\textbullet}}{\dimexpr\linewidth-26bp\relax}{نامساوی مثلثی را برای هر دو عدد مختلط ثابت کنید: \EMim{|z_1+z_2|\le|z_1|+|z_2|}.}
\EMgap
\begin{uncoverenv}<2->
\EMitem{\lr{\textbullet}}{\dimexpr\linewidth-26bp\relax}{فرم کلی نامساوی مثلثی را برای \EMim{n} عدد مختلط دلخواه ثابت کنید:}
\vspace{3bp}
\EMeq{|z_1+z_2+\cdots+z_n|\le|z_1|+|z_2|+\cdots+|z_n|}
\end{uncoverenv}
\EMgap
\begin{uncoverenv}<3->
\EMitem{\lr{\textbullet}}{\dimexpr\linewidth-26bp\relax}{ثابت کنید: \EMim{|z_1+z_2|\ge|z_1|-|z_2|}.}
\end{uncoverenv}
\end{EMbox}
\end{minipage}
\EMside{0.0bp}{\EMfigure{1.500}{m08-triangle}{نامساوی مثلثی: مثلثی با رئوس \EMim{0}، \EMim{z_1} و \EMim{z_1+z_2}}}
\end{frame}

\def\EMtitlesize{26}\def\EMtitleskip{31}
\begin{frame}{فرم قطبی اعداد مختلط}
\begin{uncoverenv}<2->
\EMeq{z =x+iy,\qquad \left.\begin{aligned}x&=|z|\cos\theta\\ y&=|z|\sin\theta\end{aligned}\right\}\; \Longrightarrow\; z =|z|\,(\cos\theta+i\sin\theta)}
\end{uncoverenv}
\EMgap
\begin{uncoverenv}<3->
\EMeq{\theta =\operatorname{Arg}(z) =\tan^{-1}\frac yx,\qquad -\pi<\theta\le\pi}
\end{uncoverenv}
\EMgap
\begin{uncoverenv}<4->
\EMeq{\arg(z) =\operatorname{Arg}(z)+2k\pi,\qquad k\in\mathbb{Z}}
\end{uncoverenv}
\end{frame}

\def\EMtitlesize{26}\def\EMtitleskip{31}
\begin{frame}{ضرب دو عدد مختلط در فرم قطبی}
\begin{uncoverenv}<2->
\EMeq{z_1 =|z_1|(\cos\theta_1+i\sin\theta_1),\qquad z_2 =|z_2|(\cos\theta_2+i\sin\theta_2)}
\end{uncoverenv}
\EMgap
\begin{uncoverenv}<3->
\EMeq{\begin{aligned}
z_1\cdot z_2&=|z_1||z_2|(\cos\theta_1+i\sin\theta_1)(\cos\theta_2+i\sin\theta_2)\\
&=|z_1||z_2|\big[(\cos\theta_1\cos\theta_2-\sin\theta_1\sin\theta_2)\\
&\qquad\qquad+\,i(\sin\theta_1\cos\theta_2+\cos\theta_1\sin\theta_2)\big]
\end{aligned}}
\end{uncoverenv}
\EMgap
\begin{EMbox}{imp}{3}{}
\begin{uncoverenv}<4->
\EMeq{z_1\cdot z_2 =|z_1||z_2|\big[\cos(\theta_1+\theta_2)+i\sin(\theta_1+\theta_2)\big]}
\end{uncoverenv}
\end{EMbox}
\begin{EMbox}{rem}{4}{}
\begin{uncoverenv}<5->
\EMp{در هنگام ضرب دو عدد مختلط، اندازه‌ها در هم ضرب و فازها با هم جمع می‌شوند.}
\end{uncoverenv}
\end{EMbox}
\begin{EMbox}{rem}{5}{}
\begin{uncoverenv}<6->
\EMp{اگر یک عدد مختلط را به توان یک عدد طبیعی برسانیم، اندازه را به توان آن عدد طبیعی رسانده، آرگومان (فاز) را در آن عدد طبیعی ضرب می‌کنیم (به روش استقراء ثابت کنید). به عبارت دیگر:}
\vspace{3bp}
\EMeq{z^k =|z|^k(\cos k\theta+i\sin k\theta)}
\end{uncoverenv}
\end{EMbox}
\end{frame}

\def\EMtitlesize{26}\def\EMtitleskip{31}
\begin{frame}{ضرب دو عدد مختلط در فرم قطبی (ادامه)}
\begin{uncoverenv}<2->
\EMeq{z^k =\big[|z|(\cos\theta+i\sin\theta)\big]^k =|z|^k(\cos\theta+i\sin\theta)^k =|z|^k(\cos k\theta+i\sin k\theta)}
\end{uncoverenv}
\EMgap
\begin{EMbox}{imp}{2}{}
\begin{uncoverenv}<3->
\EMeq{(\cos\theta+i\sin\theta)^k =\cos k\theta+i\sin k\theta}
\end{uncoverenv}
\end{EMbox}
\end{frame}

\def\EMtitlesize{26}\def\EMtitleskip{31}
\begin{frame}{ضرب دو عدد مختلط در فرم قطبی (ادامه)}
\begin{EMbox}{exe}{1}{}
\EMp{با استفاده از رابطۀ دموار \EMim{\cos3\theta} و \EMim{\cos4\theta} را بر حسب \EMim{\cos\theta} و \EMim{\sin\theta} به دست آورید.}
\end{EMbox}
\end{frame}

\def\EMtitlesize{26}\def\EMtitleskip{31}
\begin{frame}{ریشۀ \EMim{m}-ام یک عدد مختلط}
\begin{uncoverenv}<2->
\EMeq{w =\sqrt[m]{z} =z^{1/m},\qquad z =|z|(\cos\theta+i\sin\theta)}
\end{uncoverenv}
\EMgap
\begin{uncoverenv}<3->
\EMeq{z =w^m,\qquad w =|w|(\cos\varphi+i\sin\varphi)}
\end{uncoverenv}
\EMgap
\begin{uncoverenv}<4->
\EMeq{z =|z|(\cos\theta+i\sin\theta) =|w|^m(\cos m\varphi+i\sin m\varphi)}
\end{uncoverenv}
\EMgap
\begin{uncoverenv}<5->
\EMeq{|w|^m =|z|\; \Longrightarrow\;|w| =\sqrt[m]{|z|}}
\end{uncoverenv}
\EMgap
\begin{uncoverenv}<6->
\EMeq{m\varphi =\theta+2k\pi\; \Longrightarrow\;\varphi =\frac{\theta+2k\pi}{m},\qquad k =0,1,2,\dots,m-1}
\end{uncoverenv}
\EMgap
\begin{EMbox}{imp}{6}{}
\begin{uncoverenv}<7->
\EMeq{w =\sqrt[m]{z} =\sqrt[m]{|z|}\left[\cos\!\left(\frac{\theta+2k\pi}{m}\right)+i\sin\!\left(\frac{\theta+2k\pi}{m}\right)\right],\qquad k =0,1,2,\dots,m-1}
\end{uncoverenv}
\end{EMbox}
\end{frame}

\def\EMtitlesize{26}\def\EMtitleskip{31}
\begin{frame}{ریشۀ \EMim{m}-ام یک عدد مختلط (ادامه)}
\begin{minipage}[t]{500bp}\vspace{0pt}
\begin{EMbox}{ex}{1}{}
\EMp{ریشه‌های معادلۀ \EMim{z^n=1} چیست؟}
\EMgap
\begin{uncoverenv}<2->
\EMeq{z =\sqrt[n]{1},\qquad 1 =1\measuredangle0\; \Longrightarrow\; z \\[4pt] =\cos\!\left(\frac{2k\pi}{n}\right)+i\sin\!\left(\frac{2k\pi}{n}\right),\qquad k =0,1,\dots,n-1}
\end{uncoverenv}
\EMgap
\begin{uncoverenv}<3->
\EMp{با تعریف \EMim{\omega=1\measuredangle\dfrac{2\pi}{n}} ریشه‌ها عبارت‌اند از}
\vspace{3bp}
\EMeq{z =1,\ \omega,\ \omega^2,\ \dots,\ \omega^{n-1}}
\end{uncoverenv}
\end{EMbox}
\end{minipage}
\EMside{249.6bp}{\uncover<4->{\EMfigure{0.916}{m08-roots-of-unity}{ریشه‌های \EMim{z=\sqrt[3]{1}}، \EMim{z=\sqrt[4]{1}} و \EMim{z=\sqrt[5]{1}} روی دایرۀ واحد}}}
\end{frame}

\def\EMtitlesize{26}\def\EMtitleskip{31}
\begin{frame}{ریشۀ \EMim{m}-ام یک عدد مختلط (ادامه)}
\begin{EMbox}{exe}{1}{}
\EMitem{\lr{\textbullet}}{\dimexpr\linewidth-26bp\relax}{ریشه‌های معادلۀ \EMim{z^6+6+8i=0} را به دست آورده در صفحۀ اعداد مختلط رسم کنید.}
\EMgap
\begin{uncoverenv}<2->
\EMitem{\lr{\textbullet}}{\dimexpr\linewidth-26bp\relax}{درستی اتحاد \EMim{1+z+z^2+\cdots+z^n=\dfrac{1-z^{n+1}}{1-z}} را تحقیق کنید و سپس اتحاد زیر (اتحاد لاگرانژ) را نتیجه بگیرید.}
\vspace{3bp}
\EMeq{1+\cos\theta+\cos2\theta+\cdots+\cos n\theta =\frac12+\frac{\sin\!\left[\left(n+\frac12\right)\theta\right]}{2\sin\frac\theta2}}
\end{uncoverenv}
\end{EMbox}
\end{frame}
